\section{\"Uberblick}
%------------------------------------------------------------------------
\mode<presentation>{
  \begin{frame}
    \begin{enumerate}
       \item Beschreibung Fraunhofer
    \end{enumerate}
  \end{frame}
}

%------------------------------------------------------------------------
\begin{frame}[t] \frametitle{Von \TeX{} zu \LaTeX{} und \texttt{beamer}}
  \begin{description}
    \item[1978] erstes \TeX{} von Donald Knuth, vorgestellt auf 
                dem Treffen der AMS
    \item[1983] erstes \TeX{}-Manual von Leslie Lamport
    \item[1984] \LaTeX{} 2.06a von Leslie Lamport
    \item[1989] Frank Mittelbach "ubernimmt die Pflege \LaTeX{}
    \item[1994] \LaTeX{} 2.09e
    \item[2003] Till Tantau erstellt \texttt{beamer} f"ur die 
                Verteidigung seiner Doktorarbeit
    \item[2010] Joseph Wright und Vedran Mileti{\'c} "ubernehmen die 
                Pflege von \texttt{beamer}
  \end{description}

  \vskip 4mm
  \begin{columns}
    \begin{column}{0.7\textwidth}
      {\footnotesize \textsf{Quellen:}
        \begin{itemize} 
          \item \href{http://www.xent.com/FoRK-archive/feb98/0307.html}{www.xent.com/FoRK-archive/feb98/0307.html}
          \item The \texttt{beamer} class -- User guide for version 3.10
        \end{itemize}
      }
    \end{column}
  \end{columns}
\end{frame}

%------------------------------------------------------------------------
\begin{frame}[t,label=mab_refs] \frametitle{Literatur zum Kurs}
  \begin{itemize}
    \item \textsf{\LaTeX{}}:
      \begin{itemize}
        \item\label{it:SchlagerThibud} 
              Petra Schlager und Manfred Thibud.
              Wissenschaftlich mit \LaTeX arbeiten.
              Pearson Studium, M"unchen 2005.
        \item\label{it:DalheimerGuenther}
              Matthias Kalle Dalheimer und Karsten G"unther.
              \LaTeX kurz \& gut.
              3.~Auflage, O'Reilly, K"oln 2008.
      \end{itemize}
    \item \texttt{beamer}\texttt{-Klasse}:
      \begin{itemize}
        \item\label{TantauWrightMiletic}
              Till Tantau, Joseph Wright, Vedran Mileti{\'c}.
              The \texttt{beamer} class -- User guide for version 3.10.
              \href{http://www.ctan.org/tex-archive/macros/latex/contrib/beamer/doc/beameruserguide.pdf}{\url{www.ctan.org/tex-archive/macros/latex/contrib/beamer/doc/beameruserguide.pdf}}, 2010.
        \item\label{it:Voss}
              Herbert Vo"s.
              Pr"asentationen mit \LaTeX.
              Lehmanns Media (dante), Berlin 2009.
      \end{itemize}
    \item \textsf{Wissenschaftliches Arbeiten}:
      \begin{itemize}
        \item\label{it:Lagendijk}
              Ad Lagendijk.
              Survival guide for scientists -- writing -- presentation -- email.
              Amsterdam University Press, Amsterdam 2008.
      \end{itemize}
  \end{itemize}
\end{frame}

%------------------------------------------------------------------------
\begin{frame} \frametitle{Bildschirm oder nicht -- da"s ist hier die Frage}
  \mode<presentation>{
    \hskip 15.5mm \includegraphics[width=0.7\textwidth]{pics/question}
  }
  \mode<article>{
    Je nach Formeldichte, Notwendigkeit von Bildern, 
    und pers"onlichem Geschmack sollte die Wahl eines 
    Pr"asentationsmediums den Erwartungen des Publikums
    angepa"st werden. In Frage kommt etwa auch ein
    \begin{itemize}
      \item Tafelvortrag,
      \item Vortrag mit Flipchart,
      \item freier Vortrag ohne weitere Untermauerung.
    \end{itemize}
  }
\end{frame}

